\section*{Open Questions (5 new questions raised by this replication)}

\begin{description}
  \item[Q1. Non-unital noise scaling.] Our depolarizing-only replication got clean $O(\epsilon)$, $O(\epsilon^2)$, $O(\epsilon^3)$ scaling for raw, ZNE1 and ZNE2 (slopes 1.00 / 2.00 / 2.99 on $\epsilon \in [10^{-4}, 10^{-3}]$). The paper's Fig.~1(b) amplitude-damping / dephasing curves visibly flatten out at higher noise than the depolarizing curves. Does the $O(\lambda^{n+1})$ scaling law hold with the same prefactor for non-unital channels, or does the amplitude-damping steady-state bias push a channel-specific correction into the leading remainder $R_{n+1}$?
  \emph{Basis:} Only depolarizing noise was tested; Richardson's theorem is stated for any weak, time-scale-invariant $L(\rho)$, but non-unital channels behave differently under global rescaling.
  \item[Q2. Digital rescaling vs pulse-stretching.] We scaled noise by simply rebuilding a Qiskit \verb|NoiseModel| at rate $c\epsilon$. On real hardware the paper's rescaling protocol (Eq.~5) requires physically slowing the pulses by $c$ and re-calibrating gates. Does gate-duration folding (repeat each gate $c$ times) reproduce ZNE with the same order-of-cancellation as our exact digital rescaling, or does per-gate overhead / calibration drift eat one order?
  \emph{Basis:} Our slopes matched theory to 3 decimal places because our noise rescaling is exact — a luxury unavailable on hardware; Mitiq / Qiskit-Runtime use unitary-folding as a proxy whose fidelity to Eq.~5 is not benchmarked in the paper.
  \item[Q3. Optimal Richardson order $n$.] At $\epsilon = 10^{-2}$ our ZNE2 still won 8/8 vs ZNE1 but the improvement ratio (raw/ZNE2 $= 60\times$) is dramatically smaller than at $\epsilon = 10^{-3}$ (raw/ZNE2 $= 5000\times$). Where exactly is the crossover between the perturbative regime (Richardson wins big) and the regime where higher-order Born-series terms dominate? Can $n_{\text{opt}}$ be predicted analytically from ($\epsilon$, depth)?
  \emph{Basis:} The paper's bound $|E^* - \hat{E}^n_K| \le \tfrac{l_{n+1}(\lambda T)^{n+1}}{(n+1)!}\prod_j |\gamma_j|$ shows increasing $n$ eventually hurts because of the growing coefficient sum, but the crossover is not mapped in our data.
  \item[Q4. Observable-weight dependence.] Our observable was fixed at $\langle Z_0 Z_1\rangle$. The paper's Fig.~1 aggregates over ``randomly chosen multi-qubit Pauli operators $P_\alpha$''. Does the ZNE improvement factor depend on the weight/support of the observable? Concretely, do global Paulis (weight-$N$) mitigate worse than local ones because the noise couples more paths?
  \emph{Basis:} We fixed a weight-2 observable to keep the run cheap; the paper averages over random Paulis without breaking down by weight.
  \item[Q5. Shot-noise tradeoff.] The Richardson coefficients $(2,-1)$ and $(3,-3,1)$ have $L_1$-norm 3 and 7 respectively --- a variance-amplification cost invisible in our shot-free density-matrix sim. What sample-count budget $T_{\text{shots}}$ would make ZNE2 stop beating ZNE1 at each of our five noise rates?
  \emph{Basis:} Real hardware uses shots; the variance of the ZNE estimator scales like $(\sum_j |\gamma_j|)^2 / T_{\text{shots}}$, but no numerical variance curve is shown in the paper.
\end{description}
